\section{Auditable V100 Counter Assessment}
\label{sec:v100-assessment}
The campaign assigns an explicit assessment to every one of 537 names and four versions. Registry completeness does not imply completed measurements. Figure~\ref{fig:v100-coverage} exposes missing measurements rather than treating them as zeros. Each scientific-family cell counts benchmark names, so a benchmark with many launches does not dominate its domain. Source-reviewed allocation and copy cases without profiled device kernels remain outside a scalar floating-point roofline; integer, address, atomic and special-function costs are not captured by scalar add/multiply/FMA counts.

\begin{figure*}[t]\centering
\includegraphics[width=\textwidth]{figures/v100_coverage_assessment.pdf}
\caption{Dated partial-campaign assessment with the same 537-name denominator per version. Each family cell reports names with a qualified geometric position over its family denominator and separately counts unmeasured or unresolved names. Known no-counted-FP and source-reviewed inapplicable cases are separate. Samples can be partial, first-invocation or source-unassigned; this is measurement coverage, not complete elapsed-ROI coverage or causal diagnosis.}\label{fig:v100-coverage}
\Description{Four status bars distinguish qualified floating-point observations, no counted floating-point work, source-reviewed inapplicability, hardware barriers and missing measurements; a domain matrix shows qualified-name coverage and unresolved counts.}
\end{figure*}

\paragraph{Complete selected scientific families.}
A stricter source-counted audit admits twelve captures: four versions each of 100 Triad launches, ten STREAM-dot launches, and twenty Jacobi stencil launches. Exact source, configuration, result, selector and counter hashes and emitted launch counts bind the selected retained invocation. The integrated coordinate is $\sum_k F_k/\sum_k B_k$, with no equal-weight average of kernel intensities. Figure~\ref{fig:v100-integrated} reports executed FP32 instructions and measured DRAM bytes from cold kernel replay. It does not measure complete elapsed-host throughput: Jacobi's timed reset/copy, scalar transfers and host norms remain outside its stencil family, and STREAM covers dot only.

\begin{figure}[t]\centering
\includegraphics[width=\columnwidth]{figures/v100_paired_family_intensity.pdf}
\caption{Integrated executed FP32 work per measured DRAM byte for twelve source-counted scientific-family captures. Cold kernel replay is separate from the original five checked trials. This coordinate is neither an ordinary performance speedup nor a causal bottleneck label. The earlier controlled Triad useful-work intervention remains a separate experiment.}\label{fig:v100-integrated}
\Description{Triad intensities are almost equal across versions. JACC has lower integrated intensity in the selected STREAM dot and Jacobi stencil families, while native Jacobi has more counted FP32 work than the three Julia versions.}
\end{figure}

Triad's executed intensity is approximately 0.172 FLOP/DRAM byte in every version. STREAM dot is approximately 0.281 for native CUDA, CUDA.jl and KA, versus 0.0482 for JACC: the counted FP32 work is equal, while observed DRAM traffic differs. In Jacobi, native CUDA has 858.85 million counted FP32 operations versus 502.33 million in each Julia version. Intensities are 1.320, 0.770, 0.770 and 0.411, respectively. All versions perform 502.33 million counted core field-and-square operations. The additional 356.52 million native additions come from its shuffle/block norm reduction; Julia accumulates norms with conditional per-point floating atomics, whose arithmetic is excluded by the selected add/multiply/FMA counters. Thus the counted numerator also reflects reduction lowering and the metric boundary, rather than additional useful field updates or complete floating-point work. These observations identify hypotheses for additional source/SASS and intervention work. They do not establish why a version is slower, and the three families cannot support a corpus-wide intensity--speedup correlation.

\paragraph{Scientific granularity guides measurement.}
A reviewed Rayleigh--B\'enard timestep illustrates why counting authored kernels is insufficient for application interpretation. For the selected $500\times1000$ configuration, a timestep containing the Nusselt-number update has 28 authored dispatches but 12,167 source-counted cuBLAS API calls. An API call can launch zero, one or multiple device kernels; this is a source inventory, not an observed GPU launch census. The reviewed sampling proposal marks a retained-trial timestep while preserving the complete 1,000-step warmup and all five checked 1,000-step trials. Its first-per-specialization samples remain partial. Such a record connects the scientific integrator and diagnostic to backend dependence and profiling cost without assigning the whole workflow one limiting resource.

\paragraph{Observer cost and profiler method strata.}
Profiling outcomes also remain separate from translation feedback and primary performance. The Rayleigh--B\'enard observer attempt completed four successful checks of 13,720,023 values each, then reached its 1,800-second observer deadline (1,806.27 seconds including quiescence); without the fifth check, it supplies no accepted counter position. Its unchanged primary mean over trials two--five is 27.3456 seconds. Interception/replay overhead motivates investigation but has not been decomposed or established as the unique cause, and this incomplete observer run is not a primary-performance regression. Operational stops, missing payloads and parser recoveries retain their own evidence and process units; profiler versions require separately calibrated method strata rather than silently pooled measurements.

\paragraph{Source membership and sampled phase answer different questions.}
Three subsequent source-review batches bind 242 source callsite inclusion declarations to 150 selected model records across 39 names, with exact capture, configuration, result, source and stage hashes. These declarations can share a kernel family: two new BM3D call sites are aliases, and no aggregate count of independent families is inferred. The review adds no GPU executions or complete elapsed-ROI profiles; the twelve integrated scientific-family captures above remain unchanged. Scientific interpretation also depends on the actual operator. The selected all-pairs-distance contract counts integer Hamming-like symbol mismatches, rather than floating-point Euclidean distances. AMGmk's primary elapsed sum includes CPU matrix--vector and axpy work around GPU FP64 relaxation, so the GPU family cannot represent the whole sum. ASTA's in-place transpose moves FP32 payloads through integer cycle-following and atomic control; payload precision does not establish floating-point arithmetic. Barnes--Hut's timed callback contains initialization and tree construction as well as force and integration stages. A source-proved member may also execute in warmup: observed sample phase, retained invocation, launch multiplicity and whole-ROI or causal conclusions still require separate evidence.

\paragraph{Payload precision does not define the timed operation.}
A separate source review of all four versions of P2P, ping-pong and ring identifies twelve primary timed transport regions, outside a scalar compute roofline. P2P's FP32 doubling and ping-pong's FP64 increment checks occur outside these timers. The four CCL versions instead time an FP32 vendor reduction and require a compatible two-device counter method. These are source-scope findings, not new GPU observations or changes to numerical acceptance; transport latency and throughput remain meaningful performance quantities.

\paragraph{Measurement failure taxonomy and audit.}
Profiler contention, process deadlines, observer paths, host-checker dependencies, CSV parsing limits and truncated observer files are measurement failures, distinct from rejected translations and unsupported hardware. Available counter fragments, logs and captured bindings are retained; absent evidence remains explicit, and fragments from unsuccessful processes do not enter positioning. A CPU-only parser recovery retains its successful parent execution and exact original counters and counts zero new GPU or compiler attempts. Local large-file exceptions require declared payload sizes and storage reserves; a zero exit status does not certify a complete output file. Source-reviewed no-kernel applicability retains the original five-trial checks and timing. The frozen manifest audit recomputes historical derivations; it does not retroactively establish an unproved source stage or a causal resource limit. Shared review/profiling tokens remain separate from per-version translation effort.

\paragraph{Scientific requirements and implementation choices.}
The ontology links a scientific stage and accuracy contract to its reduction, precision, memory placement and backend. Table~\ref{tab:v100-ontology-interpretation} shows bounded interpretations from existing evidence. These associations help domain users choose further experiments; they do not rank abstractions for an entire scientific discipline. Source-inferred resource labels remain separate from measured counter geometry.

% Existing evidence only; no new classifications.
\begin{table}[t]\centering\small\setlength{\tabcolsep}{3pt}
\caption{Ontology fields delimit scientific interpretation. Counter observations are cold-replay selected-family measurements; source eligibility is distinct from measurement.}
\label{tab:v100-ontology-interpretation}
\begin{tabular}{p{0.23\columnwidth}p{0.34\columnwidth}p{0.32\columnwidth}}
\hline
Recorded distinction & Existing evidence & Scientific implication\\\hline
Diagnostic and reduction & Jacobi: tree adds versus floating atomics; 858.85M versus 502.33M counted FP32 operations & Same field updates need not imply the same counted-work numerator.\\
Operator and traffic & STREAM dot: equal counted work; JACC DRAM traffic 5.83 times native & Investigate this reduction implementation; no universal frontend penalty follows.\\
Memory tier & Mapped CPU-resident arrays in zerocopy (source review) & HBM-only coordinates cannot describe the primary data path.\\
Stage and backend & Scientific timesteps can contain host diagnostics and vendor calls & Authored kernels alone do not define an application's ROI.\\\hline
\end{tabular}
\end{table}

% Dated immutable partial snapshot; refresh only from a fresh audited census.
\paragraph{Operational attrition is separate from translation failure.}
At the frozen 21 October 2026, 18:38 UTC snapshot, 1,387 of 2,148 selected model slots had retained counter observations: 360 native CUDA, 344 CUDA.jl, 341 KA and 342 JACC. Every version retains the same 537-name denominator. The remaining slots include 475 pending or active measurements, 144 GPU ownership-guard stops, 26 operational timeouts, 56 plan/isolation gates, 20 model assessments sharing five hardware prerequisites, four replay-method holds and seven source-reviewed cases without profiled device kernels. Eighteen further slots have build, helper, I/O or checker diagnostics; eleven have nonzero process status with unresolved cause. These are measurement outcomes, not counts of rejected translations or unsupported framework features.

\paragraph{Selected outcomes and historical attempts have different units.}
The snapshot retains 1,732 frozen attempt records (431, 434, 434 and 433 by version). Repeated attempts and one CPU-only counter-derivation recovery are included, so these counts are not independent benchmark executions. Historical records contain 195 GPU ownership-guard stops and 33 operational timeouts, compared with 144 and 26 in the selected-slot view. Nineteen slots retain a valid earlier counter observation despite a later unsuccessful attempt. Neither the latest failure nor a recovered file changes the original five-trial timing series. A CPU-only parser recovery contributes zero new GPU or compiler executions.

\paragraph{Evidence and missingness.}
A GPU ownership-guard stop records a protective decision, not independently proved foreign contention in every case. A later assertion precheck retained three successful expected-error checks but stopped on an unrecognized process identity; its ownership could not be reconstructed after exit. That attempt remains unsuccessful, and no cause is backfilled. Historical diagnostic codes are preserved while their presentation states this uncertainty.

Categories follow a declared priority over saved stop reasons and diagnostic text. Grouped labels do not establish a unique root cause. Available counter fragments, logs and captured bindings are preserved, while absent evidence remains explicit. Forty-three historical records have no recorded profiler returncode, which can include a process that was never launched; they are not inferred missing executions. Recorded new profiler-launcher counts are 420, 425, 421 and 422, and recorded build-launcher counts are 421, 35, 34 and 35. These units may contain multiple target or compiler children and do not identify device-JIT attempts. Underlying child and device-JIT counts remain unknown. Shared measurement and paper-support usage windows are reported separately from translation windows, without adding overlapping windows or cached/reasoning subsets. The snapshot's zero manifest-integrity issues certify retained-byte joins, not full scientific, elapsed-ROI or causal coverage.

\input{data/v100_assessment/measurement_failures/selected_table}
\input{data/v100_assessment/measurement_failures/all_frozen_attempts_table}
